\chapter{Modelo de elemento de pala y cantidad de movimiento (BEMT) del rotor en autorrotación vertical}\label{cap:bemt}
\usepgfplotslibrary{groupplots}

El modelo BEMT con las polares extendidas de la placa curvada al \SI{7}{\percent}
predice, para el rotor base ($B=4$, $c=\SI{45}{mm}$, $\theta=\ang{-3}$, sin torsión) y con
el drogue atado, un descenso de $V_2\approx\SI{6.9}{m/s}$ a $\approx\num{1500}$~rpm. La
velocidad es un \SI{8}{\percent} mayor que la estimación base de \SI{6.4}{m/s}, porque el
coeficiente de resistencia del rotor que resulta es $C_R\approx0{,}97$ y no 1,2. Las rpm
son un \SI{30}{\percent} mayores que las \num{1150} de la base y confirman la estimación
por sección neutra del capítulo~\ref{cap:polares}. Este capítulo deduce el modelo,
discute el cierre de cantidad de movimiento en el estado de molino turbulento, lo valida,
barre paso, cuerda, número de palas y torsión, y analiza el arranque desde
$\Omega=0$ a \SI{13.4}{m/s}. También entrega el mapa $T(\Omega,V)$, $Q(\Omega,V)$ que usa
la simulación de vuelo.

\begin{resultado}[Resultados principales del capítulo]
\begin{itemize}
  \item Punto de diseño ($\theta=\ang{-3}$, $c=\SI{45}{mm}$, $B=4$, con drogue):
        $V_2=\SI{6.94}{m/s}$, $\Omega=\SI{157.6}{rad/s}$ (\num{1505}~rpm), punta a
        \SI{31.5}{m/s}, $T=\SI{3.58}{N}$, $C_R=0{,}97$, $V_2/v_h=2{,}03$. Sin aporte del
        drogue: $V_2=\SI{7.9}{m/s}$. Ambos valores cumplen C4 ($5\pm\SI{3}{m/s}$).
  \item $C_R$ depende sobre todo del cierre de cantidad de movimiento en el estado de
        molino turbulento. La curva de helicópteros (Castles y Gray, polinomio de
        Leishman) da $C_R=0{,}83$--$1{,}00$ y $V_2=6{,}9$--\SI{7.3}{m/s}. La corrección
        de Glauert--Buhl de turbinas eólicas da $C_R=0{,}9$--$1{,}5$ y
        $V_2=5{,}9$--\SI{7.2}{m/s}. Se adopta la primera porque fue medida en rotores en
        descenso vertical. La incertidumbre de modelo es de unos $\pm\SI{0.6}{m/s}$.
  \item Las rpm caen al subir el paso: \num{1840} a \ang{-8}, \num{1505} a \ang{-3} y
        \num{1280} a \ang{0}. En cambio, $V_2$ apenas varía ($6{,}86$--\SI{7.33}{m/s}).
        Con la sección trabajando a $C_l\approx0{,}47$ en $0{,}75R$, el valor de diseño
        $C_l\approx0{,}9$ del capítulo~\ref{cap:dim} no se alcanza.
  \item Arranque: con $\theta\le\ang{-1}$ el par a $\Omega=0$ y $V=\SI{13.4}{m/s}$ es
        positivo y $Q(\Omega)$ no tiene otra raíz que el equilibrio rápido. No hay rama
        lenta que trabe el rotor. Con $\theta\ge\ang{0}$ el rotor no arranca solo. Con
        $\theta=\ang{-3}$, $Q_0=\SI{17.8}{mN.m}$ (no los \SI{26}{mN.m} de la fórmula
        simplificada) y el rotor alcanza el \SI{90}{\percent} de su velocidad final en
        \SI{2.3}{s}, tras caer unos \SI{25}{m}.
  \item La torsión lineal no aporta: entre \ang{-8} y \ang{+8} cambia $V_2$ en menos de
        \SI{0.1}{m/s}. Con $\theta=\ang{-2}$, la torsión positiva casi anula el par
        mínimo de arranque. \textbf{Recomendación}: $\theta=\ang{-3}$ (tolerancia
        $\pm\ang{1}$), $c=\SI{45}{mm}$, $B=4$, pala sin torsión.
\end{itemize}
\end{resultado}

%=======================================================================================
\section{Formulación}\label{sec:bemt-form}

El modelo divide la pala en anillos independientes. En cada uno iguala el empuje que da
la teoría de elemento de pala con el que exige un balance de cantidad de movimiento
corregido. Se supone descenso puramente axial, flujo estacionario y anillos sin
interacción radial. Se desprecia la conicidad, cuyo efecto se estima en la
sección~\ref{sec:bemt-base}.

\subsection{Cinemática de la sección}

$V>0$ es la velocidad de descenso. El aire llega al disco desde abajo. Se define el
factor de inducción axial $a=v_i/V$, de modo que el flujo neto a través del anillo vale
$V(1-a)$. Con $a=1$ ese flujo es nulo: es la autorrotación ideal. Si el radio es $r$ y el
factor tangencial es $a'$, las componentes de la velocidad relativa son
\begin{equation}
  U_P=V(1-a),\qquad U_T=\Omega r(1+a'),\qquad \phi=\arctan\frac{U_P}{U_T},\qquad
  \alpha=\phi+\theta(r).
  \label{eq:bemt-cinem}
\end{equation}
El paso $\theta$ es negativo cuando el borde de ataque baja. Se refiere a $0{,}75R$, con
torsión lineal $\theta(r)=\theta_{0,75}+\theta_{tw}(r/R-0{,}75)$. El número de Reynolds
local es $Re=c\sqrt{U_P^2+U_T^2}/\nu$.

\subsection{Elemento de pala}

Con $C_l(\alpha,Re)$ y $C_d(\alpha,Re)$ de las polares extendidas, las fuerzas sobre el
anillo de espesor $\mathrm{d}r$ son
\begin{align}
  \mathrm{d}T &= B\,\tfrac12\rho U^2 c\,\big(C_l\cos\phi+C_d\sin\phi\big)\,\mathrm{d}r,
  \label{eq:bemt-dT}\\
  \mathrm{d}Q &= B\,\tfrac12\rho U^2 c\,\big(C_l\sin\phi-C_d\cos\phi\big)\,r\,\mathrm{d}r.
  \label{eq:bemt-dQ}
\end{align}
El empuje $T$ es positivo hacia arriba. El par $Q$ es positivo cuando impulsa el giro.
Las ecuaciones coinciden con~\eqref{eq:dT} y~\eqref{eq:dQ} del capítulo de teoría. Si
$\mathrm{d}Q>0$ el anillo es \emph{impulsor}. Si $\mathrm{d}Q<0$ es \emph{impulsado}. Si
$\alpha$ supera el de $C_{l,\max}$ ($\approx\ang{7.5}$) está \emph{en pérdida}.

Dividiendo~\eqref{eq:bemt-dT} por $\tfrac12\rho V^2\,2\pi r\,\mathrm{d}r$ se obtiene el
coeficiente de empuje local de elemento de pala, con la solidez local
$\sigma_r=Bc/(2\pi r)$:
\begin{equation}
  C_{T,\mathrm{EP}}=\sigma_r\,\frac{U^2}{V^2}\,\big(C_l\cos\phi+C_d\sin\phi\big).
  \label{eq:bemt-ctbe}
\end{equation}

\subsection{Pérdidas de punta y raíz}

El número finito de palas se tiene en cuenta con el factor de Prandtl~\cite{leishman2006}.
Se aplica en la punta y en la raíz, porque la raíz está en $e=\SI{44}{mm}$ y allí también
se desprende un torbellino:
\begin{equation}
  F=\frac{4}{\pi^2}\arccos\!\Big[e^{-\frac{B}{2}\frac{R-r}{r\sin\phi}}\Big]\,
         \arccos\!\Big[e^{-\frac{B}{2}\frac{r-e}{r\sin\phi}}\Big].
  \label{eq:bemt-F}
\end{equation}
Como $\phi$ es pequeño (2--\ang{8}), $F$ cae con rapidez solo en el último
\SI{3}{\percent} del radio.

\subsection{Cierre de cantidad de movimiento: el problema del molino turbulento}\label{sec:bemt-cierre}

En el estado de molino frenante ($a<0{,}5$) el balance de cantidad de movimiento da
$C_T=4aF(1-a)$. Esa relación pierde validez para $a>0{,}5$, porque predice una estela
que invierte su sentido. El rotor del CanSat trabaja en $a\approx0{,}7$--$0{,}86$, es
decir, en el estado de \emph{molino turbulento}, en la región $-2<V_c/v_h<0$ donde la
teoría de cantidad de movimiento no vale (fig.~\ref{fig:bemt-curvas}). Hay dos curvas
empíricas para sustituirla.

\begin{itemize}
\item \textbf{Helicópteros}: Castles y Gray~\cite{castles1951} midieron rotores en
  descenso vertical en túnel y Leishman~\cite{leishman2006} ajustó sus datos con el
  polinomio~\eqref{eq:poly}. Con $x=V/v_h$ y $s=V/v_h$ en el descenso, ese polinomio da
  $v_i/v_h=f(-s)$ y, por lo tanto,
  \begin{equation}
    a=\frac{f(-s)}{s},\qquad C_T=\frac{T}{\tfrac12\rho V^2A}=\frac{4}{s^2}.
    \label{eq:bemt-heli}
  \end{equation}
  La curva paramétrica $(a(s),C_T(s))$ se une a la rama $4a(1-a)$ en el punto donde ambas
  se cortan, $a_x=\num{0.39}$ ($s=2{,}07$). Así queda una función continua y monótona
  $C_T(a)$, que se multiplica por $F$. En $a=1$ da $C_T=1{,}21$, el $C_R$ de la
  autorrotación ideal del capítulo~\ref{cap:teoria}.
\item \textbf{Turbinas eólicas}: la corrección de Glauert, en la forma de
  Buhl~\cite{bemt-buhl2005}, es una parábola que empalma con $4aF(1-a)$ en $a=0{,}4$ y
  vale $C_T=2$ en $a=1$:
  \begin{equation}
    C_T=\frac89+\Big(4F-\frac{40}{9}\Big)a+\Big(\frac{50}{9}-4F\Big)a^2 \quad (a>0{,}4).
    \label{eq:bemt-buhl}
  \end{equation}
\end{itemize}

\begin{figure}[H]
\centering
\begin{tikzpicture}
\begin{axis}[width=0.8\textwidth, height=6.6cm, xlabel={Factor de inducción $a$},
  ylabel={$C_T$ local}, xmin=0, xmax=1.4, ymin=0, ymax=2.6, grid=major,
  legend pos=north west, legend style={font=\footnotesize}]
\addplot[cgris, dashed, thick] table[x=a, y=CTmom] {analysis/data/bemt_curvas.dat};
\addlegendentry{$4a(1-a)$ (cantidad de movimiento)}
\addplot[cfreno, thick] table[x=a, y=CTbuhl] {analysis/data/bemt_curvas.dat};
\addlegendentry{Glauert--Buhl, $F=1$}
\addplot[cfreno, thin, densely dotted] table[x=a, y=CTbuhl09] {analysis/data/bemt_curvas.dat};
\addlegendentry{Glauert--Buhl, $F=0{,}9$}
\addplot[cfijo, very thick] table[x=a, y=CTheli] {analysis/data/bemt_curvas.dat};
\addlegendentry{Helicópteros (Leishman), $F=1$}
\addplot[cfijo, thin, densely dotted] table[x=a, y=CTheli09] {analysis/data/bemt_curvas.dat};
\addlegendentry{Helicópteros, $F=0{,}9$}
\fill[crot, opacity=0.15] (axis cs:0.68,0) rectangle (axis cs:0.87,2.6);
\node[font=\footnotesize, crot!80!black, align=center] at (axis cs:0.775,0.25) {zona de\\trabajo};
\end{axis}
\end{tikzpicture}
\caption{Cierres de cantidad de movimiento en el estado de molino turbulento. La banda
naranja marca el intervalo de $a$ en que trabajan los anillos del rotor en el punto de
diseño.}
\label{fig:bemt-curvas}
\end{figure}

En la zona de trabajo ($a\approx0{,}75$) las dos curvas difieren en un
\SI{25}{\percent}: helicópteros da $C_T\approx0{,}9$ y Glauert--Buhl, $\approx1{,}2$. La
diferencia tiene un origen físico. La curva de Glauert se ajustó con datos de turbinas y
hélices muy cargadas, con solidez y regímenes que no son los de un rotor en descenso
vertical. Además, la forma de Buhl se diseñó solo para que fuera continua con la rama de
molino frenante, no para reproducir $C_T(a=1)$. Los datos de Castles y Gray, en cambio,
son de rotores en descenso vertical libre de túnel: es justamente la situación del
CanSat. Por eso el modelo base usa la curva de helicópteros. Glauert--Buhl se usa como
cota del lado favorable (descenso más lento), y la incertidumbre del propio ajuste se
representa escalando un $\pm\SI{20}{\percent}$ el exceso de $C_T$ sobre la rama de
cantidad de movimiento (sec.~\ref{sec:bemt-sens}).

\begin{nota}[Límite de validez]
Ninguna de las dos curvas describe un flujo estacionario real. En el molino turbulento
la estela recircula y es inestable, y el BEMT solo da valores medios. En la rama
$a>1$ (rotor casi parado en el arranque) la curva de helicópteros se prolonga como
$C_T\propto a^2$, que corresponde a una placa porosa. Por eso los pares del arranque
(sec.~\ref{sec:bemt-arranque}) son estimaciones de orden de magnitud.
\end{nota}

\subsection{Algoritmo}\label{sec:bemt-alg}

Para cada par $(\Omega,V)$ se resuelve el residuo de cada anillo,
\begin{equation}
  g(a)=C_{T,\mathrm{EP}}(a)-C_T^{\mathrm{mom}}(a;F),
  \label{eq:bemt-res}
\end{equation}
con este procedimiento:
\begin{enumerate}
  \item Arranque desde $a=0{,}5$. Iteración de punto fijo
        $\hat a=\big(C_T^{\mathrm{mom}}\big)^{-1}\!\big(C_{T,\mathrm{EP}}(a)\big)$, con
        relajación dinámica de Aitken por anillo, $a\leftarrow a+w(\hat a-a)$,
        $w\in[10^{-3},1]$.
  \item Salvaguarda: $g$ decrece con $a$, de modo que se mantiene un intervalo
        $[a_\mathrm{lo},a_\mathrm{hi}]$ que contiene la raíz. Si el paso sale del
        intervalo o el residuo no baja al menos un \SI{30}{\percent}, se da un paso de
        bisección.
  \item Tolerancia $|g|<10^{-9}$. Si no se alcanza en 200 iteraciones, el anillo se
        resuelve por bisección pura (55 pasos).
  \item Con todos los anillos convergidos se integran~\eqref{eq:bemt-dT}
        y~\eqref{eq:bemt-dQ} por la regla del punto medio, con 40 anillos uniformes
        entre $e$ y $R$.
\end{enumerate}
La inducción tangencial $a'$ es de orden $10^{-3}$ porque el par neto es casi nulo. Se
toma $a'=0$ y su efecto se verifica en la sección~\ref{sec:bemt-sens}.

\paragraph{Equilibrio.} El equilibrio vertical exige dos condiciones:
\begin{equation}
  Q(\Omega,V)=0,\quad \frac{\partial Q}{\partial\Omega}<0,\qquad
  T(\Omega,V)+\tfrac12\rho V^2\big(C_{D,d}A_d+C_{D,c}A_c\big)=W.
  \label{eq:bemt-eq}
\end{equation}
Para cada $V$ se buscan todas las raíces de $Q(\Omega)$ con un barrido de 41 puntos
(0 y una malla geométrica de 2 a \SI{500}{rad/s}) más un refinamiento por Brent, y se
toma la raíz estable más rápida. Después se resuelve en $V$ por Brent entre
\SI{2.5}{m/s} y \SI{16}{m/s}. Un equilibrio completo cuesta unas 500 evaluaciones BEMT
(unos \SI{2}{s}).

\paragraph{Polares.} Se usan las tablas extendidas de la placa al \SI{7}{\percent}
($N_\mathrm{crit}=9$, paso \ang{0.5}) para $Re=\num{20000}$, \num{35000}, \num{50000},
\num{75000} y \num{100000}. La interpolación es lineal en $\alpha$ y en $\ln Re$, y se
congela fuera de ese rango.

%=======================================================================================
\section{Validación y convergencia}\label{sec:bemt-valid}

El código reproduce la solución exacta de un caso analítico con un error de
$3\times10^{-10}$ en $a$, y el empuje converge al \SI{0.2}{\percent} con 40 anillos.

\paragraph{Caso 1: molino frenante con polar lineal.} Se toma $B=2$, $c=\SI{20}{mm}$,
$\theta=\ang{-4}$, $C_l=2\pi\alpha$, $C_d=0$, sin pérdidas, con cierre $4a(1-a)$, a
$V=\SI{8}{m/s}$ y $\Omega=\SI{250}{rad/s}$. La solución de pequeños ángulos es
\begin{equation}
  \lambda=\tfrac12\Big[\big(\lambda_\infty-\tfrac{\sigma a_0}{8}\big)
  +\sqrt{\big(\lambda_\infty-\tfrac{\sigma a_0}{8}\big)^2-\tfrac{\sigma a_0\theta x}{2}}\Big],
  \qquad a=1-\frac{\lambda}{\lambda_\infty},
  \label{eq:bemt-anal}
\end{equation}
con $\lambda_\infty=V/(\Omega R)$. Además se calculó una solución exacta escalar,
independiente de la rutina vectorial. La figura~\ref{fig:bemt-valid} muestra el
resultado. El BEMT coincide con la solución exacta a $3\times10^{-10}$ y da
$T=\SI{2.9394}{N}$. La fórmula~\eqref{eq:bemt-anal} da \SI{2.9381}{N}: difiere en un
\SI{0.05}{\percent} y en $a$ difiere en 0,006 en la raíz, donde $\phi$ llega a \ang{27}
y la aproximación de pequeños ángulos falla.

\paragraph{Caso 2: autorrotación ideal.} Con $C_d=0$, sin pérdidas y con la curva de
helicópteros, la autorrotación debe dar $a\approx1$ en todo el disco y
$C_R\to C_T(a=1)=1{,}21$. El código da $a=0{,}87$--1,08 y $C_R=1{,}15$. La diferencia
del \SI{5}{\percent} se debe a que, con paso constante, el ángulo de flujo no es
uniforme: la raíz impulsa con $a<1$ y la punta frena con $a>1$.

\begin{figure}[H]
\centering
\begin{tikzpicture}
\begin{groupplot}[group style={group size=2 by 1, horizontal sep=1.8cm},
  width=0.47\textwidth, height=5.4cm, grid=major, xlabel={$r/R$},
  legend style={font=\scriptsize}, legend pos=north west]
\nextgroupplot[ylabel={$a$}, title={\footnotesize Caso 1: molino frenante}, ymin=0.12, ymax=0.32,
  legend pos=north east]
\addplot[cfijo, thick] table[x=x, y=a_exac] {analysis/data/bemt_valid_lineal.dat};
\addlegendentry{exacta}
\addplot[only marks, mark=o, mark size=1.3pt, crot] table[x=x, y=a_bemt] {analysis/data/bemt_valid_lineal.dat};
\addlegendentry{BEMT}
\addplot[cfreno, dashed] table[x=x, y=a_anal] {analysis/data/bemt_valid_lineal.dat};
\addlegendentry{ec.~\eqref{eq:bemt-anal}}
\nextgroupplot[ylabel={$a$}, title={\footnotesize Caso 2: autorrotación ideal}, ymin=0.8, ymax=1.12]
\addplot[crot, thick, mark=*, mark size=0.8pt] table[x=x, y=a] {analysis/data/bemt_valid_ideal.dat};
\addlegendentry{BEMT}
\addplot[cgris, dashed, domain=0.2:1] {1};
\addlegendentry{$a=1$}
\end{groupplot}
\end{tikzpicture}
\caption{Validación del código BEMT con dos casos de referencia.}
\label{fig:bemt-valid}
\end{figure}

\paragraph{Convergencia.} En el punto de diseño se repitió el cálculo con 5 a 320 anillos
(tabla~\ref{tab:bemt-malla}). Con 40 anillos el empuje está a \SI{0.23}{\percent} del
valor con 320. El par residual de \SI{-0.15}{mN.m} desplaza la $\Omega$ de equilibrio en
$\Delta\Omega=-Q/(\partial Q/\partial\Omega)\approx\SI{0.17}{rad/s}$, que es despreciable.
La iteración converge en 25 pasos con el cierre de helicópteros y en 20 con Buhl, en los
40 anillos sin recurrir a la bisección de respaldo. Una bisección independiente coincide
a $2\times10^{-10}$. Un barrido fino de $a\in[-0{,}9;\,3{,}9]$ confirmó que cada anillo
tiene una sola raíz.

\begin{table}[H]
\centering
\caption{Convergencia con el número de anillos en el punto de diseño
($\Omega=\SI{157.6}{rad/s}$, $V=\SI{6.94}{m/s}$).}
\label{tab:bemt-malla}
\begin{tabular}{rccc}
\toprule
Anillos $N$ & $T$ [N] & $Q$ [mN·m] & Error de $T$ respecto de $N=320$ \\
\midrule
5   & 3,678 & 2,377  & \SI{2.95}{\percent} \\
10  & 3,623 & 0,975  & \SI{1.42}{\percent} \\
20  & 3,595 & 0,258  & \SI{0.62}{\percent} \\
40  & 3,581 & 0,000  & \SI{0.23}{\percent} \\
80  & 3,575 & $-0{,}104$ & \SI{0.08}{\percent} \\
160 & 3,573 & $-0{,}139$ & \SI{0.02}{\percent} \\
320 & 3,573 & $-0{,}152$ & --- \\
\bottomrule
\end{tabular}
\end{table}

%=======================================================================================
\section{Punto de diseño}\label{sec:bemt-base}

Con $\theta=\ang{-3}$ el rotor desciende a \SI{6.94}{m/s} y gira a \num{1505}~rpm. El
empuje es \SI{3.58}{N}. El drogue aporta \SI{1.16}{N} y el cuerpo \SI{0.17}{N}, y entre
los tres equilibran $W=\SI{4.905}{N}$.

\subsection{Distribuciones radiales}

La figura~\ref{fig:bemt-dist} muestra el estado de cada anillo. El \SI{77}{\percent}
interior del radio ($0{,}22<r/R<0{,}78$) es impulsor y el \SI{22}{\percent} exterior es
impulsado. No hay anillos en pérdida: $\alpha$ va de \ang{6.3} en la raíz a
\ang{-0.6} en la punta, por debajo del $\alpha$ de $C_{l,\max}$. El factor $a$ va de 0,86
en la raíz a 0,71 en la punta. $Re$ va de \num{21000} a \num{94000}, y en $0{,}75R$ vale
\num{72000}, igual que en la base. En $0{,}75R$ la sección trabaja a
$\alpha=\ang{1.5}$ y $C_l=0{,}47$, dentro del intervalo 0,42--0,53 del
capítulo~\ref{cap:polares}. El $C_l$ medio ponderado con el empuje es 0,57.

\begin{figure}[H]
\centering
\begin{tikzpicture}
\begin{groupplot}[group style={group size=2 by 2, horizontal sep=1.7cm, vertical sep=1.6cm},
  width=0.47\textwidth, height=5cm, grid=major, xmin=0.2, xmax=1,
  legend style={font=\scriptsize}]
\nextgroupplot[ylabel={Ángulos [\si{\degree}]}, legend pos=north east]
\addplot[cfijo, thick] table[x=x, y=alpha] {analysis/data/bemt_dist_heli.dat};
\addlegendentry{$\alpha$}
\addplot[crot, thick, dashed] table[x=x, y=phi] {analysis/data/bemt_dist_heli.dat};
\addlegendentry{$\phi$}
\nextgroupplot[ylabel={$C_l$, $a$}, legend pos=north east, ymin=0, ymax=1.4]
\addplot[cfijo, thick] table[x=x, y=cl] {analysis/data/bemt_dist_heli.dat};
\addlegendentry{$C_l$}
\addplot[ccuerda, thick, dashed] table[x=x, y=a] {analysis/data/bemt_dist_heli.dat};
\addlegendentry{$a$}
\nextgroupplot[xlabel={$r/R$}, ylabel={$\mathrm{d}T/\mathrm{d}r$ [N/m]}, legend pos=north west]
\addplot[cfijo, thick] table[x=x, y=dT] {analysis/data/bemt_dist_heli.dat};
\addlegendentry{helicópteros}
\addplot[cfreno, dashed] table[x=x, y=dT] {analysis/data/bemt_dist_buhl.dat};
\addlegendentry{Buhl}
\nextgroupplot[xlabel={$r/R$}, ylabel={$\mathrm{d}Q/\mathrm{d}r$ [mN·m/m]}, legend pos=south west, ymin=-300, ymax=160]
\fill[ccuerda, opacity=0.12] (axis cs:0.2,-300) rectangle (axis cs:0.775,160);
\fill[cfreno, opacity=0.10] (axis cs:0.775,-300) rectangle (axis cs:1,160);
\addplot[cfijo, thick] table[x=x, y=dQmili] {analysis/data/bemt_dist_heli.dat};
\addlegendentry{helicópteros}
\addplot[cfreno, dashed] table[x=x, y=dQmili] {analysis/data/bemt_dist_buhl.dat};
\addlegendentry{Buhl}
\addplot[black, thin, domain=0.2:1] {0};
\node[font=\scriptsize, ccuerda!60!black] at (axis cs:0.5,-40) {impulsora};
\node[font=\scriptsize, cfreno] at (axis cs:0.89,100) {impulsada};
\end{groupplot}
\end{tikzpicture}
\caption{Distribuciones radiales en el punto de diseño ($\theta=\ang{-3}$, 80 anillos).
El área de $\mathrm{d}Q/\mathrm{d}r$ es nula: equilibrio de par.}
\label{fig:bemt-dist}
\end{figure}

$\mathrm{d}T/\mathrm{d}r$ crece casi linealmente hacia la punta y cae en el último
\SI{5}{\percent} por la pérdida de Prandtl. El par impulsor máximo está en
$r/R\approx0{,}45$; la punta, con $U_T$ grande y $\phi$ pequeño, es la que más frena. Con Buhl la zona
impulsora se acorta hasta $r/R=0{,}75$ y el empuje sube en toda la pala.

\subsection{Comparación con la estimación base}

\begin{table}[H]
\centering
\caption{Punto de diseño: BEMT frente a la estimación base (cap.~\ref{cap:dim}).}
\label{tab:bemt-base}
\begin{tabular}{lccc}
\toprule
Magnitud & Base & BEMT (helicópteros) & BEMT (Buhl) \\
\midrule
$V_2$ con drogue [m/s]   & 6,4  & 6,94 & 6,29 \\
$V_2$ sin drogue [m/s]   & 7,8 ($C_R=1$) & 7,91 & --- \\
$\Omega$ [rpm]           & \num{1150} & \num{1505} & \num{1548} \\
Velocidad de punta [m/s] & 24 & 31,5 & 32,4 \\
$T$ [N]                  & 3,78 & 3,58 & 3,82 \\
$C_R$                    & 1,2 & 0,97 & 1,25 \\
$V_2/v_h$                & 1,83 & 2,03 & 1,79 \\
$C_l$ en $0{,}75R$       & 0,9 & 0,47 & 0,46 \\
$Re$ en $0{,}75R$        & \num{72000} & \num{72000} & \num{73000}\\
\bottomrule
\end{tabular}
\end{table}

La tabla~\ref{tab:bemt-base} resume las diferencias, que tienen dos causas:
\begin{enumerate}
  \item \textbf{$C_R$ menor.} La base supone $C_R=1{,}2$, el valor de la autorrotación
        ideal. El BEMT da 0,97, por dos motivos: la resistencia de perfil obliga a que
        parte del disco trabaje con $a<1$ para impulsar, y $F$ reduce el empuje en la
        punta y en la raíz. Con menos empuje del rotor, el drogue tiene que aportar más,
        y para eso la velocidad sube a \SI{6.94}{m/s}. Con Buhl, $C_R=1{,}25$ y el
        resultado queda cerca de la base.
  \item \textbf{$C_l$ de trabajo menor.} La ec.~\eqref{eq:omega} supuso $C_l=0{,}9$. El
        equilibrio de par fija $\alpha$ cerca del ángulo de máxima eficiencia de
        autorrotación y no del de $C_l$ alto: con $\theta=\ang{-3}$ resulta $C_l=0{,}47$
        en $0{,}75R$. Para generar el mismo empuje con $C_l$ menor hacen falta más rpm:
        $\Omega\propto C_l^{-1/2}$ da $1150\sqrt{0{,}9/0{,}57}\approx\num{1450}$~rpm,
        cerca de las \num{1505} del BEMT.
\end{enumerate}
Esto confirma la estimación de \num{1500}--\num{1700}~rpm del
capítulo~\ref{cap:polares} y corrige las \num{1150}~rpm base. Las cargas centrífugas
crecen con $\Omega^2$, un factor 1,71: unos \SI{22}{N} por pala en lugar de \SI{13}{N}.

La conicidad de \ang{6.1} del capítulo de estructura inclina el empuje de cada pala.
Su componente vertical baja en $\cos\beta_0=0{,}994$ y el radio efectivo en el mismo
factor. Al primer orden ($T\propto R^3\cos\beta_0$), eso equivale a perder un
\SI{2}{\percent} de empuje a igual $\Omega$. Como el equilibrio se reajusta, $V_2$ sube
unos \SI{0.05}{m/s}, menos que la incertidumbre de modelo.

%=======================================================================================
\section{Barrido de paso, cuerda y número de palas}\label{sec:bemt-barrido}

El paso controla las rpm y el arranque, pero casi no influye en la velocidad de
descenso. Entre \ang{-8} y \ang{+2}, $V_2$ solo varía entre \SI{6.86}{m/s} y
\SI{7.33}{m/s} (figuras~\ref{fig:bemt-barrido-c} y~\ref{fig:bemt-barrido-B}).

\begin{figure}[H]
\centering
\begin{tikzpicture}
\begin{groupplot}[group style={group size=2 by 1, horizontal sep=1.7cm},
  width=0.47\textwidth, height=5.6cm, grid=major, xlabel={$\theta_{0,75}$ [\si{\degree}]},
  legend style={font=\scriptsize}, unbounded coords=jump, xmin=-8, xmax=2]
\nextgroupplot[ylabel={$V_2$ [m/s]}, legend pos=north east, ymin=6.7, ymax=7.6]
\addplot[cfijo, thick, mark=square*, mark size=1.2pt] table[x=theta, y=V_c35] {analysis/data/bemt_theta_c.dat};
\addlegendentry{$c=\SI{35}{mm}$}
\addplot[crot, very thick, mark=*, mark size=1.2pt] table[x=theta, y=V_c45] {analysis/data/bemt_theta_c.dat};
\addlegendentry{$c=\SI{45}{mm}$}
\addplot[ccuerda, thick, mark=triangle*, mark size=1.5pt] table[x=theta, y=V_c55] {analysis/data/bemt_theta_c.dat};
\addlegendentry{$c=\SI{55}{mm}$}
\nextgroupplot[ylabel={$n$ [rpm]}, legend pos=north east, ymin=900, ymax=2200]
\addplot[cfijo, thick, mark=square*, mark size=1.2pt] table[x=theta, y=rpm_c35] {analysis/data/bemt_theta_c.dat};
\addplot[crot, very thick, mark=*, mark size=1.2pt] table[x=theta, y=rpm_c45] {analysis/data/bemt_theta_c.dat};
\addplot[ccuerda, thick, mark=triangle*, mark size=1.5pt] table[x=theta, y=rpm_c55] {analysis/data/bemt_theta_c.dat};
\addplot[cgris, dashed, domain=-8:2] {1150};
\node[font=\scriptsize, cgris, anchor=south west] at (axis cs:-7.8,1150) {base 1150};
\end{groupplot}
\end{tikzpicture}
\caption{Efecto del paso y de la cuerda ($B=4$, con drogue).}
\label{fig:bemt-barrido-c}
\end{figure}

\begin{figure}[H]
\centering
\begin{tikzpicture}
\begin{groupplot}[group style={group size=2 by 1, horizontal sep=1.7cm},
  width=0.47\textwidth, height=5.6cm, grid=major, xlabel={$\theta_{0,75}$ [\si{\degree}]},
  legend style={font=\scriptsize}, unbounded coords=jump, xmin=-8, xmax=2]
\nextgroupplot[ylabel={$n$ [rpm]}, legend pos=north east, ymin=900, ymax=2200]
\addplot[cfijo, thick, mark=square*, mark size=1.2pt] table[x=theta, y=rpm_B3] {analysis/data/bemt_theta_B.dat};
\addlegendentry{$B=3$}
\addplot[crot, very thick, mark=*, mark size=1.2pt] table[x=theta, y=rpm_B4] {analysis/data/bemt_theta_B.dat};
\addlegendentry{$B=4$}
\addplot[ccuerda, thick, mark=triangle*, mark size=1.5pt] table[x=theta, y=rpm_B5] {analysis/data/bemt_theta_B.dat};
\addlegendentry{$B=5$}
\nextgroupplot[ylabel={$Q_0$ a \SI{13.4}{m/s} [mN·m]}, legend pos=north east]
\addplot[cfijo, thick, mark=square*, mark size=1.2pt] table[x=theta, y=Q0_B3] {analysis/data/bemt_theta_B.dat};
\addplot[crot, very thick, mark=*, mark size=1.2pt] table[x=theta, y=Q0_B4] {analysis/data/bemt_theta_B.dat};
\addplot[ccuerda, thick, mark=triangle*, mark size=1.5pt] table[x=theta, y=Q0_B5] {analysis/data/bemt_theta_B.dat};
\addplot[black, thin, domain=-8:2] {0};
\end{groupplot}
\end{tikzpicture}
\caption{Efecto del número de palas ($c=\SI{45}{mm}$): rpm de equilibrio y par de
arranque a $\Omega=0$.}
\label{fig:bemt-barrido-B}
\end{figure}

El barrido ($11\times3\times3=99$ configuraciones) muestra cuatro tendencias:
\begin{itemize}
  \item \textbf{Paso.} Cada grado de paso más negativo suma unas 70~rpm y unos
        \SI{6}{mN.m} de par de arranque. También baja el $C_l$ de trabajo y, con él,
        $C_R$ (0,83 a \ang{-8}), por lo que $V_2$ sube \SI{0.4}{m/s} entre \ang{-3} y
        \ang{-8}. El mínimo de $V_2$ ($\approx\SI{6.86}{m/s}$) está en $\theta\approx\ang{0}$,
        pero ese paso no arranca.
  \item \textbf{Solidez.} Más cuerda o más palas bajan las rpm casi como
        $(Bc)^{-1/2}$ y apenas mejoran $V_2$: \SI{0.08}{m/s} de \SI{35}{mm} a
        \SI{55}{mm}. El $Re$ en $0{,}75R$ depende de $c\,\Omega$ y cambia poco.
  \item \textbf{Par de arranque.} $Q_0$ crece con $Bc$ y con $-\theta$. A \ang{0} es
        nulo para todas las configuraciones: con la polar extendida la pala a
        $\alpha=\ang{90}$ se comporta como placa plana simétrica.
  \item \textbf{Velocidad de punta.} Con $c=\SI{35}{mm}$ y $B=3$ supera los
        \SI{40}{m/s} (\num{1960}~rpm), y con eso suben las cargas centrífugas y el ruido.
        El rotor base queda en el centro del intervalo.
\end{itemize}

\begin{table}[H]
\centering
\caption{Resultados del barrido a $\theta=\ang{-3}$ (con drogue).}
\label{tab:bemt-barrido}
\begin{tabular}{ccccccccc}
\toprule
$c$ [mm] & $B$ & $\sigma$ & $V_2$ [m/s] & rpm & $C_R$ & $C_l(0{,}75R)$ & $V_\text{punta}$ [m/s] & $Q_0$ [mN·m] \\
\midrule
35 & 3 & 0,167 & 7,06 & \num{1960} & 0,92 & 0,48 & 41,0 & 11,2 \\
35 & 4 & 0,223 & 6,99 & \num{1710} & 0,95 & 0,47 & 35,9 & 14,7 \\
35 & 5 & 0,279 & 6,94 & \num{1540} & 0,97 & 0,46 & 32,2 & 18,0 \\
45 & 3 & 0,215 & 7,00 & \num{1710} & 0,94 & 0,49 & 35,9 & 13,7 \\
\textbf{45} & \textbf{4} & \textbf{0,286} & \textbf{6,94} & \textbf{\num{1505}} & \textbf{0,97} & \textbf{0,47} & \textbf{31,5} & \textbf{17,8} \\
45 & 5 & 0,358 & 6,90 & \num{1360} & 0,98 & 0,46 & 28,5 & 21,8 \\
55 & 3 & 0,263 & 6,96 & \num{1540} & 0,96 & 0,49 & 32,3 & 15,9 \\
55 & 4 & 0,350 & 6,90 & \num{1350} & 0,98 & 0,48 & 28,3 & 20,7 \\
55 & 5 & 0,438 & 6,87 & \num{1230} & 0,99 & 0,46 & 25,7 & 25,2 \\
\bottomrule
\end{tabular}
\end{table}

Sin el drogue (rotor solo con el cuerpo), $V_2$ sube a 7,8--\SI{8.5}{m/s} según el paso
(\SI{7.91}{m/s} a \ang{-3}, con \num{1700}~rpm). Es el peor caso si el drogue se enreda
o se suelta, y cumple C4 ($\le\SI{8}{m/s}$) con $\theta\ge\ang{-3}$, pero sin margen: con
$\theta\le\ang{-4}$ ya lo supera (\SI{8.03}{m/s} a \ang{-4}).

%=======================================================================================
\section{Sensibilidad al modelo}\label{sec:bemt-sens}

El cierre del molino turbulento domina la incertidumbre de $V_2$: con Buhl, $V_2$ baja
entre \SI{0.6}{m/s} y \SI{0.9}{m/s}. Las demás hipótesis cambian $V_2$ en menos de
\SI{0.1}{m/s} y las rpm en menos de un \SI{6}{\percent} (tabla~\ref{tab:bemt-sens}).

\begin{table}[H]
\centering
\caption{Sensibilidad del punto $\theta=\ang{-3}$ a las hipótesis del modelo.}
\label{tab:bemt-sens}
\begin{tabular}{lccc}
\toprule
Variante & $V_2$ [m/s] & rpm & $C_R$ \\
\midrule
Base: helicópteros, $F$ punta y raíz, polares $Re$ local, $a'=0$ & 6,94 & \num{1505} & 0,97 \\
Cierre Glauert--Buhl                       & 6,29 & \num{1548} & 1,25 \\
Exceso de $C_T$ de helicópteros $\times0{,}8$ & 7,00 & \num{1500} & 0,94 \\
Exceso de $C_T$ de helicópteros $\times1{,}2$ & 6,87 & \num{1510} & 0,99 \\
Sin pérdidas de Prandtl ($F=1$)            & 6,85 & \num{1520} & 1,00 \\
Con inducción tangencial $a'$              & 6,94 & \num{1500} & 0,97 \\
Polar única de $Re=\num{50000}$            & 6,98 & \num{1510} & 0,95 \\
$C_d\times1{,}2$ (rugosidad, bordes)       & 6,98 & \num{1430} & 0,95 \\
Polar 2D sin corrección rotacional         & 6,95 & \num{1500} & 0,96 \\
\bottomrule
\end{tabular}
\end{table}

\begin{figure}[H]
\centering
\begin{tikzpicture}
\begin{groupplot}[group style={group size=2 by 1, horizontal sep=1.7cm},
  width=0.47\textwidth, height=5.4cm, grid=major, xlabel={$\theta_{0,75}$ [\si{\degree}]},
  legend style={font=\scriptsize}, unbounded coords=jump, xmin=-8, xmax=2]
\nextgroupplot[ylabel={$V_2$ [m/s]}, legend pos=south west, ymin=5.6, ymax=7.6]
\addplot[crot, very thick, mark=*, mark size=1.2pt] table[x=theta, y=V_base] {analysis/data/bemt_modelos.dat};
\addlegendentry{helicópteros}
\addplot[cfreno, thick, dashed, mark=o, mark size=1.2pt] table[x=theta, y=V_buhl] {analysis/data/bemt_modelos.dat};
\addlegendentry{Glauert--Buhl}
\addplot[cfijo, thick, mark=square, mark size=1.2pt] table[x=theta, y=V_cd12] {analysis/data/bemt_modelos.dat};
\addlegendentry{$C_d\times1{,}2$}
\addplot[cgris, dashed, domain=-8:2] {6.4};
\node[font=\scriptsize, cgris, anchor=south east] at (axis cs:1.8,6.4) {base 6,4};
\nextgroupplot[ylabel={$C_R$}, legend pos=north west, ymin=0.7, ymax=1.6]
\addplot[crot, very thick, mark=*, mark size=1.2pt] table[x=theta, y=CR_base] {analysis/data/bemt_modelos.dat};
\addplot[cfreno, thick, dashed, mark=o, mark size=1.2pt] table[x=theta, y=CR_buhl] {analysis/data/bemt_modelos.dat};
\addplot[cfijo, thick, mark=square, mark size=1.2pt] table[x=theta, y=CR_cd12] {analysis/data/bemt_modelos.dat};
\addplot[cgris, dashed, domain=-8:2] {1.2};
\end{groupplot}
\end{tikzpicture}
\caption{Sensibilidad de $V_2$ y $C_R$ al cierre de cantidad de movimiento y a la
resistencia de perfil.}
\label{fig:bemt-sens}
\end{figure}

Con Buhl el $C_R$ crece al subir el paso, porque los anillos se acercan a $a=1$, donde
esa curva da $C_T=2$. Con helicópteros el $C_R$ se satura cerca de 1,0. Los ensayos
publicados de autogiros y de rotores en descenso vertical dan $C_R\approx1{,}0$--1,3, y
el polinomio de Leishman en autorrotación ideal da 1,2. En ese intervalo, Buhl es
plausible a pasos moderados, pero sus valores de 1,4--1,5 a $\theta\ge\ang{0}$ ya son
altos. Como cifra de diseño se adopta $V_2=\SI{6.9}{m/s}$, con un intervalo de
\SI{6.3}{m/s} a \SI{7.0}{m/s} con drogue. El ensayo de caída del capítulo de ensayos
(medición de $V_2$ y rpm) es lo que debe decidir entre los dos cierres.

%=======================================================================================
\section{Torsión de la pala}\label{sec:bemt-torsion}

Una torsión lineal de \ang{-8} a \ang{+8} cambia $V_2$ en menos de \SI{0.1}{m/s} y las
rpm en menos de un \SI{2}{\percent}. En la práctica, la torsión no mejora el descenso.

\begin{figure}[H]
\centering
\begin{tikzpicture}
\begin{groupplot}[group style={group size=2 by 1, horizontal sep=1.7cm},
  width=0.47\textwidth, height=5.4cm, grid=major, xlabel={Torsión $\theta_{tw}$ [\si{\degree}]},
  legend style={font=\scriptsize}, xmin=-8, xmax=8]
\nextgroupplot[ylabel={$V_2$ [m/s]}, legend pos=north east, ymin=6.75, ymax=7.1]
\addplot[cfijo, thick, mark=square*, mark size=1.2pt] table[x=tw, y=V_m4] {analysis/data/bemt_torsion.dat};
\addlegendentry{$\theta=\ang{-4}$}
\addplot[crot, very thick, mark=*, mark size=1.2pt] table[x=tw, y=V_m3] {analysis/data/bemt_torsion.dat};
\addlegendentry{$\theta=\ang{-3}$}
\addplot[ccuerda, thick, mark=triangle*, mark size=1.5pt] table[x=tw, y=V_m2] {analysis/data/bemt_torsion.dat};
\addlegendentry{$\theta=\ang{-2}$}
\nextgroupplot[ylabel={$Q_{\min}$ a \SI{13.4}{m/s} [mN·m]}, ymin=0, ymax=28, legend pos=south west]
\addplot[cfijo, thick, mark=square*, mark size=1.2pt] table[x=tw, y=Qmin_m4] {analysis/data/bemt_torsion.dat};
\addplot[crot, very thick, mark=*, mark size=1.2pt] table[x=tw, y=Qmin_m3] {analysis/data/bemt_torsion.dat};
\addplot[ccuerda, thick, mark=triangle*, mark size=1.5pt] table[x=tw, y=Qmin_m2] {analysis/data/bemt_torsion.dat};
\end{groupplot}
\end{tikzpicture}
\caption{Efecto de la torsión lineal (paso referido a $0{,}75R$; $\theta_{tw}>0$ sube el
paso de la punta). A la derecha, el par mínimo de $Q(\Omega)$ entre 0 y el equilibrio
durante el arranque a \SI{13.4}{m/s}.}
\label{fig:bemt-torsion}
\end{figure}

La torsión negativa (punta con menos paso) da más par en la raíz, pero también aumenta
la zona impulsada. La positiva sube un poco $C_R$ (0,967 a 0,985 con \ang{+8} a
\ang{-3}). El efecto importante está en el arranque. Con $\theta=\ang{-2}$ y torsión
positiva, la raíz queda con paso casi nulo o positivo, y el par mínimo durante la
aceleración baja de \SI{11.9}{mN.m} a \SI{0.25}{mN.m} con \ang{+2}. El rotor quedaría al
borde de trabarse a pocas rpm. Una pala plana es además más fácil de fabricar, así que
se recomienda \textbf{no torcer la pala}.

%=======================================================================================
\section{Arranque y equilibrios múltiples}\label{sec:bemt-arranque}

Con $\theta=\ang{-3}$ el rotor arranca solo y sin riesgo de quedar trabado. En todo el
intervalo $V=2$--\SI{16}{m/s}, $Q(\Omega)$ es positivo desde $\Omega=0$ y tiene una
única raíz, que es estable.

\begin{figure}[H]
\centering
\begin{tikzpicture}
\begin{groupplot}[group style={group size=2 by 1, horizontal sep=1.7cm},
  width=0.47\textwidth, height=5.8cm, grid=major, xlabel={$\Omega$ [rad/s]},
  legend style={font=\scriptsize}, xmin=0, xmax=400]
\nextgroupplot[ylabel={$Q$ [mN·m]}, title={\footnotesize $\theta=\ang{-3}$, varias $V$},
  ymin=-30, ymax=30, legend pos=south west, legend columns=2]
\addplot[cfijo!40, thick] table[x=Om, y=Q_V4p0] {analysis/data/bemt_qomega_V.dat};
\addlegendentry{4}
\addplot[cfijo!70, thick] table[x=Om, y=Q_V6p0] {analysis/data/bemt_qomega_V.dat};
\addlegendentry{6}
\addplot[cfijo, thick] table[x=Om, y=Q_V8p0] {analysis/data/bemt_qomega_V.dat};
\addlegendentry{8}
\addplot[ccont, thick] table[x=Om, y=Q_V10p0] {analysis/data/bemt_qomega_V.dat};
\addlegendentry{10}
\addplot[crot, very thick] table[x=Om, y=Q_V13p4] {analysis/data/bemt_qomega_V.dat};
\addlegendentry{13,4}
\addplot[cfreno, thick] table[x=Om, y=Q_V16p0] {analysis/data/bemt_qomega_V.dat};
\addlegendentry{16 m/s}
\addplot[black, thin, domain=0:400] {0};
\nextgroupplot[ylabel={$Q$ [mN·m]}, title={\footnotesize $V=\SI{13.4}{m/s}$, varios $\theta$},
  ymin=-30, ymax=50, legend pos=north east, legend columns=2]
\addplot[cfijo, thick] table[x=Om, y=Q_tm8] {analysis/data/bemt_qomega_th.dat};
\addlegendentry{$-8$}
\addplot[cfijo!60, thick] table[x=Om, y=Q_tm6] {analysis/data/bemt_qomega_th.dat};
\addlegendentry{$-6$}
\addplot[crot, very thick] table[x=Om, y=Q_tm3] {analysis/data/bemt_qomega_th.dat};
\addlegendentry{$-3$}
\addplot[ccuerda, thick] table[x=Om, y=Q_tm1] {analysis/data/bemt_qomega_th.dat};
\addlegendentry{$-1$}
\addplot[cgris, thick, dashed] table[x=Om, y=Q_t0] {analysis/data/bemt_qomega_th.dat};
\addlegendentry{$0$}
\addplot[cfreno, thick] table[x=Om, y=Q_t2] {analysis/data/bemt_qomega_th.dat};
\addlegendentry{$+2$}
\addplot[black, thin, domain=0:400] {0};
\end{groupplot}
\end{tikzpicture}
\caption{Curvas de par $Q(\Omega)$. Izquierda: paso recomendado a distintas
velocidades. Derecha: distintos pasos a la velocidad de la etapa 1.}
\label{fig:bemt-qomega}
\end{figure}

La forma de $Q(\Omega)$ explica el comportamiento (fig.~\ref{fig:bemt-qomega}). Con
$\Omega=0$ la pala ve el aire a $\alpha=\ang{90}+\theta$. Con $\theta<0$ la componente
de la fuerza normal en el plano empuja la pala hacia adelante: es el par de arranque.
Al subir $\Omega$, $\alpha$ baja por la zona de pérdida profunda, donde el par se
mantiene positivo y decrece en forma monótona hasta el equilibrio. Con
$\theta\le\ang{-1}$ no aparece ningún mínimo local que cruce el cero. Con
$\theta=\ang{+1}$, en cambio, $Q(0)<0$ y aparece una raíz inestable a 9--\SI{51}{rad/s}
($V=3$--\SI{16}{m/s}). El rotor solo llega a la rama rápida si algo lo lleva por encima
de esa raíz. Si arranca parado, queda frenado en $\Omega=0$
(fig.~\ref{fig:bemt-locus}). Esto confirma, con el modelo completo, la regla del
capítulo~\ref{cap:dim}: con paso $\ge\ang{0}$ el rotor no arranca solo.

\begin{figure}[H]
\centering
\begin{tikzpicture}
\begin{axis}[width=0.75\textwidth, height=6cm, grid=major, xlabel={$V$ [m/s]},
  ylabel={$\Omega$ con $Q=0$ [rad/s]}, xmin=3, xmax=16, ymin=0, ymax=420,
  legend pos=north west, legend style={font=\scriptsize}, unbounded coords=jump]
\addplot[cfijo, thick] table[x=V, y=Om_est] {analysis/data/bemt_locus_tm6.dat};
\addlegendentry{$\theta=\ang{-6}$, estable}
\addplot[crot, very thick] table[x=V, y=Om_est] {analysis/data/bemt_locus_tm3.dat};
\addlegendentry{$\theta=\ang{-3}$, estable}
\addplot[ccuerda, thick] table[x=V, y=Om_est] {analysis/data/bemt_locus_tm1.dat};
\addlegendentry{$\theta=\ang{-1}$, estable}
\addplot[cfreno, thick] table[x=V, y=Om_est] {analysis/data/bemt_locus_t1.dat};
\addlegendentry{$\theta=\ang{+1}$, estable}
\addplot[cfreno, thick, dashed] table[x=V, y=Om_inest] {analysis/data/bemt_locus_t1.dat};
\addlegendentry{$\theta=\ang{+1}$, inestable (umbral)}
\addplot[only marks, mark=*, mark size=2pt, black] coordinates {(6.94,157.6)};
\node[font=\scriptsize, anchor=west] at (axis cs:7.1,150) {diseño};
\end{axis}
\end{tikzpicture}
\caption{Lugar de equilibrio de par $Q(\Omega,V)=0$. Con $\theta=\ang{+1}$ aparece una
raíz inestable a bajas rpm que hace de umbral: por debajo de ella el rotor se frena.}
\label{fig:bemt-locus}
\end{figure}

\paragraph{Par a $\Omega=0$.} El BEMT da $Q_0=\SI{17.8}{mN.m}$ con \ang{-3} y
\SI{11.9}{mN.m} con \ang{-2}. Son menos de la mitad del resultado de la fórmula
simplificada~\eqref{eq:Q0} (\SI{26}{mN.m} con \ang{-2}). La fórmula usa
$C_N=2\sin\alpha$ y la velocidad libre completa. El BEMT, en cambio, usa la polar
extendida a \ang{90} ($C_N=1{,}17$, corregido por alargamiento, frente a 1,98 en 2D) e incluye el
bloqueo del disco, que con el rotor parado frena el flujo ($a>1$ en la rama
extrapolada). Con la polar 2D, sin corrección de alargamiento, $Q_0$ sube a
\SI{25.7}{mN.m}. El arranque solo exige que $Q_0$ supere el par de rozamiento de los rodamientos y
los sellos, que no se midió; conviene medirlo en el banco (estimación: pocas décimas
de mN·m a \SI{1}{mN.m}).

\paragraph{Transitorio de arranque.} Con el mapa de la sección siguiente se integró el
arranque cuasiestacionario:
\begin{equation}
  I_z\dot\Omega=Q(\Omega,V),\qquad m\dot V=W-T(\Omega,V)-\tfrac12\rho V^2(C_{D,d}A_d+C_{D,c}A_c),
  \label{eq:bemt-trans}
\end{equation}
con $I_z=\SI{5.1e-4}{kg.m^2}$ y condiciones iniciales $\Omega=0$, $V=\SI{13.4}{m/s}$. La
figura~\ref{fig:bemt-trans} muestra la evolución. Durante el primer segundo la aceleración
es lenta, porque el par es pequeño. Entre 1,5 y \SI{2.5}{s} la aceleración es fuerte.
$\Omega$ supera el valor final en un \SI{4.5}{\percent}, con un máximo de
\SI{164.8}{rad/s} a los \SI{2.95}{s}, y se estabiliza en \SI{157.8}{rad/s} y
\SI{6.94}{m/s}. El \SI{90}{\percent} de $\Omega$ final se alcanza a los \SI{2.3}{s},
tras caer \SI{25}{m}. A los \SI{6}{s} la caída acumulada es de \SI{52}{m}. La dinámica
del batimiento, la apertura de las palas y la inercia del flujo (que el modelo
cuasiestacionario ignora) se tratan en el capítulo de dinámica.

\begin{figure}[H]
\centering
\begin{tikzpicture}
\begin{axis}[width=0.75\textwidth, height=5.8cm, grid=major, xlabel={$t$ [s]},
  ylabel={$n$ [rpm]}, xmin=0, xmax=6, ymin=0, ymax=1800, axis y line*=left,
  legend style={font=\scriptsize}]
\addplot[crot, very thick] table[x=t, y=rpm] {analysis/data/bemt_transitorio.dat};
\label{pl:bemt-rpm}
\end{axis}
\begin{axis}[width=0.75\textwidth, height=5.8cm, xmin=0, xmax=6, ymin=0, ymax=15,
  axis y line*=right, axis x line=none, ylabel={$V$ [m/s]},
  legend pos=north east, legend style={font=\scriptsize}]
\addlegendimage{/pgfplots/refstyle=pl:bemt-rpm}\addlegendentry{rpm}
\addplot[cfijo, thick, dashed] table[x=t, y=V] {analysis/data/bemt_transitorio.dat};
\addlegendentry{$V$}
\end{axis}
\end{tikzpicture}
\caption{Arranque cuasiestacionario desde $\Omega=0$ y $V=\SI{13.4}{m/s}$
($\theta=\ang{-3}$, con drogue).}
\label{fig:bemt-trans}
\end{figure}

%=======================================================================================
\section{Mapa para la simulación}\label{sec:bemt-mapa}

El archivo \texttt{analysis/data/bemt\_mapa.dat} contiene $T$ [N] y $Q$ [N·m] para
$\theta=\ang{-3}$, $c=\SI{45}{mm}$, $B=4$, sin torsión, con el cierre de helicópteros. La
rejilla es regular: $\Omega=0$--\SI{200}{rad/s} cada \SI{5}{rad/s} y $V=2$--\SI{16}{m/s}
cada \SI{0.5}{m/s} (\num{1189} puntos). Las columnas son \texttt{Omega V T Q}. La
interpolación bilineal es suficiente, porque $T$ y $Q$ son suaves en toda la rejilla (no
hay saltos de rama). Fuera de la rejilla conviene extrapolar con $T\propto V^2$ y
$Q\propto V^2$ a $\Omega/V$ constante, que es la similitud del molino.

%=======================================================================================
\section{Conclusiones y recomendación}\label{sec:bemt-concl}

\begin{resultado}[Configuración recomendada]
\begin{itemize}
  \item \textbf{Paso $\theta_{0,75}=\ang{-3}$}, aceptable entre \ang{-2} y \ang{-4}.
        Con \ang{-2} el par de arranque baja a \SI{12}{mN.m}. Con \ang{-4} las rpm suben
        a \num{1580} y $V_2$ a \SI{7.0}{m/s}. Nunca $\ge\ang{0}$.
  \item \textbf{Cuerda \SI{45}{mm} y $B=4$}: se mantienen. Más solidez casi no mejora
        $V_2$ y agrega masa. Menos solidez sube la velocidad de punta por encima de
        \SI{35}{m/s}.
  \item \textbf{Sin torsión}: la ganancia es menor que \SI{0.1}{m/s}, y la torsión
        positiva compromete el arranque.
  \item \textbf{Prestaciones esperadas}: $V_2=\SI{6.9}{m/s}$ (intervalo de modelo
        6,3--\SI{7.0}{m/s}) con drogue y \SI{7.9}{m/s} sin drogue, \num{1500}~rpm
        (\num{1430}--\num{1550}) y arranque en $\approx\SI{2.3}{s}$.
  \item \textbf{Cambios respecto de la base}: $V_2$ es \SI{0.5}{m/s} mayor
        ($C_R=0{,}97$ en lugar de 1,2), las rpm son un \SI{30}{\percent} mayores
        (\num{1505} en lugar de \num{1150}), la fuerza centrífuga sube a unos \SI{22}{N}
        por pala y el par de arranque es de \SI{18}{mN.m} con \ang{-3} (no
        \SI{26}{mN.m} con \ang{-2}). Los capítulos de estructura y de dinámica deben usar
        estos valores.
\end{itemize}
\end{resultado}

Los límites del modelo son tres: el cierre empírico del molino turbulento, la hipótesis
de flujo estacionario y anillos independientes, y las polares calculadas con XFOIL (no
medidas). La validación definitiva debe hacerse con el ensayo de caída instrumentado,
midiendo rpm con el sensor hall y $V_2$ con el barómetro.
